\subsection*{Impactos sociales}
\addcontentsline{toc}{subsection}{Impactos sociales}

El despliegue de la plataforma proyecta una incidencia significativa en el tejido social de Bogotá, orientada principalmente a la democratización del conocimiento financiero. A continuación, se describen las dimensiones de este impacto:

\begin{itemize}
	\item \textbf{Inclusión financiera y reducción de la brecha educativa:} La plataforma actúa como un agente de alfabetización digital y financiera, abordando directamente la carencia de competencias técnicas en segmentos tradicionalmente excluidos por el sistema bancario formal. Al proporcionar herramientas didácticas tales como simuladores de tasas de interés y motores de cálculo de amortización, se empodera al usuario para transitar de una inclusión financiera nominal a una participación activa y crítica en el mercado de capitales.
	
	\item \textbf{Mitigación del estrés financiero y salud mental:} El bienestar financiero se define como el estado de seguridad que permite cumplir obligaciones económicas presentes y futuras. La incapacidad de gestionar deudas o presupuestos genera patologías asociadas a la ansiedad financiera que afectan la productividad laboral y la estabilidad personal. Al proveer mecanismos de seguimiento de gastos y establecimiento de metas de ahorro, el aplicativo reduce la incertidumbre económica, promoviendo una mejora en la calidad de vida y la salud mental de los usuarios.
	
	\item \textbf{Fortalecimiento del capital humano:} El acceso a contenido educativo personalizado mediante IA permite que los jóvenes y trabajadores independientes desarrollen habilidades de planificación estratégica. Esto no solo mejora su capacidad de ahorro, sino que optimiza la toma de decisiones sobre inversión en educación, emprendimiento y protección social, factores determinantes para la movilidad socioeconómica en el contexto colombiano.
	
	\item \textbf{Efecto multiplicador y transformación cultural:} A largo plazo, se prevé un impacto intergeneracional basado en la transferencia de hábitos financieros saludables. La adopción de metodologías de presupuesto y ahorro programado dentro del núcleo familiar fomenta una cultura de responsabilidad financiera que rompe con ciclos de endeudamiento crónico, generando un beneficio social que trasciende al usuario individual.
\end{itemize}